\section{Related Work}
\label{sec:related-work}

\subsection{Eye Movements as a Reading Signal}
Decades of reading research establish that oculomotor behaviour reflects cognitive processing: fixation durations track processing effort, saccades relocate attention, and regressions, backward saccades to earlier text, are associated with comprehension difficulty~\cite{rayner1998eyemovements, rayner2009eyemovements}.
Eye-movement measures predict reading performance in online reading contexts~\cite{guan2022predictive}, and recent work shows that global measures such as fixation counts and regression rates predict scores on standardised comprehension tests~\cite{meziere2023predicting}.
Salvucci and Goldberg's taxonomy of fixation-identification algorithms~\cite{salvucci2000identifying} remains the reference point for converting raw gaze into oculomotor events; our platform implements a dwell-based, area-of-interest variant over live text tokens.
This literature justifies gaze as the primary sensing channel for reading difficulty, but it does not by itself provide the infrastructure to act on the signal in real time.

\subsection{Gaze-Contingent and Adaptive Reading Interfaces}
Typography measurably affects reading for both general and impaired readers: letter spacing~\cite{beier2021letterspacing}, font size~\cite{rello2016makeitbig}, background colour~\cite{rello2017backgroundcolors}, and reader-controlled text customisation~\cite{arditi2004adjustable} all shift performance or comfort.
Systems research has closed the loop between gaze and text since eyeBook's gaze-responsive reading experiences and the reading assistance of Beymer et al.~\cite{beymer2008eyetracking}.
Gaze Tutor detected disengagement from gaze and triggered tutorial dialogue moves~\cite{dmello2012gazetutor}; the Eye-Mind Reader detected mind wandering during reading and intervened with self-explanation prompts, improving long-term comprehension~\cite{hutt2020eyemind}, and the approach has since been carried into real classrooms~\cite{hutt2021classrooms}.
Within the Reading the Reader project, three predecessor prototypes demonstrated gaze-driven typographic adaptation for readers with central vision loss~\cite{refsgaard2024rtr, pereira2024typography, palinec2024reactive}, and Jensen et al.\ introduced reading resume time as a measure of post-intervention disruption, finding that only a gradual intervention paired with highlighting significantly reduced recovery time~\cite{jensen2025context}.
Micro-intervention design frameworks likewise argue for small, just-in-time adaptations over wholesale changes~\cite{persson2025microintervention}.
Across this line of work the adaptive loop is repeatedly rebuilt: each system binds one sensor, one detection method, and one intervention repertoire into one interface, so intervention comparisons across systems remain confounded by the platforms themselves.

\subsection{Experiment Tooling}
General-purpose experiment software approaches the problem from the opposite side.
Tobii Pro Lab~\cite{tobii2025prolab} designs and records eye-tracking studies but does not drive real-time text adaptation; PsychoPy~\cite{peirce2019psychopy} and Psychtoolbox offer stimulus control and gaze recording with scripting, but adaptive reading interfaces, pluggable decision modules, and researcher-in-the-loop intervention control must be built from scratch.
Consumer reading tools such as Immersive Reader~\cite{microsoft2025immersive} adapt presentation but sense nothing.
To our knowledge, no available tool combines experiment design, gaze recording, real-time gaze-driven text intervention, and pluggable decision modules; this capability gap is what our platform targets.

\subsection{Toward Low-Cost Sensing}
Research-grade trackers such as the Tobii Pro Fusion~\cite{tobii2025fusion} deliver the sampling rates and accuracy that reading research needs, but their cost and setup confine studies to the laboratory.
Webcam-based eye tracking has been validated as a preliminary instrument for reading tests~\cite{lin2022webcam}, and commodity cameras can additionally supply facial signals that research trackers do not observe at all.
A platform whose sensing layer is a replaceable module can adopt such sensors as they mature without rebuilding the experiment infrastructure around them; Section~\ref{sec:webcam} reports our work in progress in this direction.
